\section*{Chapter \ref{ch:iv2:behavioural-and-fit} --- Behavioural and Fit}
\markboth{\protect\omTocMark{}{Answers: Behavioural and Fit}}{\protect\omTocMark{}{Answers: Behavioural and Fit}}

\begin{solution}{iq:iv2:behavioural-and-fit:1}
Give a short story that ends at this job: what you studied or did, the one or two things you did that show you can do this work, and why you are here today. Stop at about ninety seconds, and let them choose what to ask about.

To build your own, write three sentences of past, two of evidence and one of present. Choose evidence with a number in it, that you are ready to be questioned on for ten minutes.

\textbf{Illustration, with your own facts in the brackets.} ``I studied [subject] at [place], where I found I liked [specific kind of problem]. I then worked on [project], where I [what you did], and the result was [number]. Separately, I built [second thing], which taught me [skill]. I am applying for this role because [specific reason tied to the work]. I am happy to go deeper on any of that.''

\iqwrong{Reading out your CV in order, or starting with your childhood. It takes five minutes, and it tells them nothing they cannot already read.}
\end{solution}

\begin{solution}{iq:iv2:behavioural-and-fit:2}
Name something specific that this work has and the other does not: for example, that the result of your work is measured quickly and in money, that your code or your model is close to a decision, or that the problem itself interests you. Do not say that it is about pay, and do not pretend pay does not matter.

To build your own, list what you have enjoyed most in your actual work, and find the feature that finance adds to it. Then think about what you give up: products that many people see and use, work you can publish, and the freedom to share what you learn. Say that you know.

\textbf{Illustration, with your own facts in the brackets.} ``In [project], what I liked most was [getting a fast answer on whether I was right]. In finance, that loop is short and hard to argue with, because [your reason]. I understand you often cannot publish, and that a good idea stops working. I would take that for [your reason].''

\iqwrong{``I have always been passionate about markets,'' with nothing behind it. They will ask what you have read, traded or built. Have an answer ready.}
\iqfollow{Is it the money?}{Be honest: pay matters, and it is not the only reason. A flat denial is less believable, and it starts you off badly.}
\end{solution}

\begin{solution}{iq:iv2:behavioural-and-fit:3}
Say which part of the work you want to do all day, and give evidence that you already do it by choice. Show that you respect the other role, and do not describe it as second best.

Roughly, the difference to discuss is speed against depth. Trading rewards fast decisions under uncertainty, with a result you see the same day. Research rewards long, careful work where most ideas fail and the result arrives after months. Which you prefer should be visible in what you have done.

\textbf{Illustration for a trader, with your own facts in the brackets.} ``When I [played a game or ran a position, in your words], I liked deciding with incomplete information and seeing the result quickly. On [long project], I found I was less patient at the [late stage]. So I would rather be closest to the decision, and to work with researchers who have more patience for the proof than I do.''

\iqwrong{``I'm not sure, I would do either.'' It sounds as if you have not looked into either. If you truly have no preference, say what would decide it, and that you want to find out by doing a specific thing here.}
\end{solution}

\begin{solution}{iq:iv2:behavioural-and-fit:4}
Compare what the work is like, rather than ranking the firms: how money is made, how close you are to it, how quickly you learn whether you were right, and what the risks are. Then give the reason that suits you.

To build your own, prepare one sentence for each type, in terms of what you have read, and check each one against public sources. Roughly: a market maker earns many small amounts by quoting prices and managing the risk, a hedge fund holds views for longer, a bank serves clients and manages regulated risk, and a crypto firm works in younger markets with more operational risk. Treat that as a starting point, because firms differ widely.

\textbf{Illustration, with your own facts in the brackets.} ``I want [type] because I want [feature], and I saw that when I [your evidence]. I looked at [other type], and I would miss [feature]. That is a preference, not a judgement on them.''

\iqwrong{Criticising the other three. You may be talking to someone who used to work at one, and it says more about you than about them.}
\end{solution}

\begin{solution}{iq:iv2:behavioural-and-fit:5}
Give a short, specific account from public sources: what the firm trades, in which markets, how it earns, and what could hurt it. Then say what you do not know, and ask.

To prepare, read the firm's own pages, its published work and talks, any public filings and recent news. Write down three facts, one question and one risk. The risk should be real: for a firm that quotes prices, for example, volumes falling, competitors getting faster, or a technical failure, each of which can be argued from how such a firm earns.

\textbf{Illustration, with your own facts in the brackets.} ``From [public source], I understand you mainly do [activity] in [markets], so you earn when [mechanism]. I would think your main risks are [risk one] and [risk two]. I could not find out [specific thing]. Is that something you can tell me about?''

\iqwrong{Repeating the home page, or naming a rumour from a forum as fact. Rumours about pay and methods are often out of date and sometimes false. If you do cite one, say where it came from and that you cannot vouch for it.}
\end{solution}

\begin{solution}{iq:iv2:behavioural-and-fit:6}
Name a real weakness that matters a little for this job, give an example, and say what you do about it and how far that has worked.

To find one, ask yourself what a colleague or teacher last told you to improve, or what went wrong in the project from your worst week. Do not choose something fatal to the role. A trader who says they freeze under pressure has answered the wrong question.

\textbf{Illustration, with your own facts in the brackets.} ``I tend to [real habit, such as polishing too long before showing work]. On [project], it cost [concrete effect]. Now I [specific change, such as showing a rough version by day two], and that has helped, although I still [what remains].''

\iqwrong{``I'm a perfectionist,'' or ``I work too hard.'' Everyone hears these as dodges. It also tells them you won't give a straight answer when you find something awkward, which on a desk is the weakness that matters most.}
\end{solution}

\begin{solution}{iq:iv2:behavioural-and-fit:7}
Describe the work you want to be doing and how good you want to be at it. Do not give a title, or a plan to leave. Add what you would have learned by then, and say that you may be wrong about the details.

To build your own, find out what someone five years in actually does at this kind of firm, and ask them at the interview if you can. Your plan should show that you want depth here. The time scale of your plan should match the role's.

\textbf{Illustration, with your own facts in the brackets.} ``I would like to be someone who owns [specific kind of work] and who can [skill you lack today]. To get there, I expect to spend the first years learning [what]. I do not know yet whether that means [route one] or [route two], and that is one of the things I want to find out here.''

\iqwrong{``Running my own fund,'' or ``in your seat.'' Ambition is fine, but this tells them that you see the role as a stepping stone. An answer of ``I haven't thought about it'' is just as weak.}
\end{solution}

\begin{solution}{iq:iv2:behavioural-and-fit:8}
Do not compare yourself with people you have never met. Give the two or three things that you bring, each with evidence, and tie them to what this role needs.

To build your own, read the role's description and pick its hardest part. Then choose something that you have done which most people with your background have not: for example, a thing that you built, a result you can show, or an unusual combination of skills. Say it plainly, and then stop. A good answer takes less than a minute.

\textbf{Illustration, with your own facts in the brackets.} ``I can't speak about the others. What I can tell you is that I have [evidence one with a number] and [evidence two]. For this role, that means I could [concrete contribution] sooner than someone who would need to learn it. And I work well with [kind of people], for example when [short true example].''

\iqwrong{Being so modest that you offer nothing, or claiming to be the best. They asked you to make a case, and neither of those is one.}
\end{solution}

\begin{solution}{iq:iv2:behavioural-and-fit:9}
Here is a one-minute version you can adapt. ``Think of a shop that buys and sells used cars. It will buy your car at 9\,000 and sell one like it at 9\,500, and it earns the 500 by doing that all day, for anyone who comes in. It does not need to know whether car prices will rise. A market maker does the same with shares and other investments, but on screens, in tiny amounts, many times a second. It loses money in two ways. If prices fall while it holds a lot of stock, it is stuck with the loss. And if the person selling knows something that it does not, such as a problem with the car, it has paid too much. So it keeps its stock small, and moves its prices quickly.''

To build your own, choose one image the listener already knows, and test it on someone who is not in finance.

\iqwrong{Using the words spread, liquidity, inventory or adverse selection. If you need those words, you have not explained it. Do not forget the part on losing money, which is what the question is really about.}
\iqfollow{So why do people trade with it, if it always takes a cut?}{Because they can trade at once, whenever they want, for a small cost. Without it they might wait a long time to find a buyer.}
\end{solution}

\begin{solution}{iq:iv2:behavioural-and-fit:10}
Say what happened in one or two plain sentences, take responsibility for your part, say what changed, and give evidence that it changed. Then stop. Do not apologise repeatedly, and do not tell a long story.

To build your own, decide what is true and how much you want to share. You need not give private details, such as a health or family matter: ``a family health matter, now resolved'' is enough. Your evidence should come after the weak spot: a later result, a finished project.

\textbf{Illustration, with your own facts in the brackets.} ``In [year], my results fell to [result]. The reason was [short true reason]. I changed [what you did], and after that I [later result]. It is a fair question, and I think [later result] is the better guide to what I can do.''

\iqwrong{Blaming a teacher, a manager or a course, or changing the dates on your CV. A check can reveal that, and it would end the process. Hiding bad news is what worries an employer most.}
\iqfollow{How do I know it won't happen again?}{Say what is different now, and give a time that you were under the same pressure afterwards.}
\end{solution}

\begin{solution}{iq:iv2:behavioural-and-fit:11}
Be honest that you are interviewing elsewhere, say broadly where, and say what would decide your choice. You do not have to give names, and you should never invent an offer.

If you have a deadline, say so early and calmly: ``I have an offer that expires on [date], so it would help me to know where I stand by then.'' That is usually better for both sides than a surprise at the last moment.

\textbf{Illustration, with your own facts in the brackets.} ``Yes, I am talking to [number] other firms of [kind]. What will decide it for me is [what you care about]. From what I have seen, this firm is strong on [true thing], and I am still finding out about [other thing].''

\iqwrong{``You are my only choice.'' It is rarely believed, and if it is untrue then your first statement to them is a lie. The other error is playing firms against each other with an offer you do not have. That is a small world, and people talk.}
\end{solution}

\begin{solution}{iq:iv2:behavioural-and-fit:12}
Say what you actually do, and what it did and did not teach you. Games can teach you about thinking in odds, sizing and losing calmly. They do not show that you can trade. If you do not play, say so, and give what you do have.

To build your own, be ready for a number: how you know whether you win. A person who plays for money but cannot say their long-run result, or how they recorded it, has not learned the main lesson.

\textbf{Illustration, with your own facts in the brackets.} ``I play [game] [how often], and I kept records over [period], which showed [result]. It taught me [lesson about sizing or tilt, in your words]. Where it differs from trading is that [real difference].'' If you do not play: ``I don't play. Where I have thought hardest about odds and bad results is [your example].''

\iqwrong{Telling stories about one big win, or implying that a gambling habit is the same thing as trading. Someone who bets without an edge, or to feel something, is exactly what a risk manager fears.}
\end{solution}

\begin{solution}{iq:iv2:behavioural-and-fit:13}
Choose something hard and not on your syllabus. Say how you picked your sources, how you practised, and what test showed you had learned it. The last part is what most people leave out.

To build your own, describe the difficulty honestly: what you got wrong first, and how you found that out. ``I read a book'' is not an account of learning. A test might be a working program, a solved set of problems, an explanation that someone expert accepted, or a result checked against a known answer.

\textbf{Illustration, with your own facts in the brackets.} ``I taught myself [topic] over [time], because [reason]. I started with [source], and I got stuck on [point]. So I [what you did]. I knew it had worked when [test, such as: my code matched a published number, or I solved a set of problems with no help].''

\iqwrong{Choosing something easy, or something from a course, because it feels safe. Choosing something you cannot discuss in depth is worse, because they may know it better than you do.}
\end{solution}

\begin{solution}{iq:iv2:behavioural-and-fit:14}
Ask two or three real questions that you could not have answered from public sources, about the person's own work and how success is judged. Keep questions about pay and terms for after an offer.

For a future manager, ask what a new person gets wrong in the first six months, how you would be judged after a year, what the team is working on that is hard, and how disagreements get settled. For a senior partner, ask what has changed in the business in the last few years, what worries them, and what they look for in people who stay. After an offer, ask about pay and how it is set, notice periods and any restriction on working elsewhere after you leave, how much of the code or research you can use elsewhere, and how a position or project moves to someone else if it fails.

\iqwrong{``No, you have covered everything.'' It suggests little interest. Asking something on the website, or asking first about holidays or bonus size, is also a poor signal.}
\iqfollow{What if they answer something vaguely?}{Ask one follow-up, politely, and note it. A vague answer on how people are judged is information you need for your decision.}
\end{solution}

\begin{solution}{iq:iv2:behavioural-and-fit:15}
Choose a real disappointment that you owned and can explain with numbers. State what you expected and why, what happened, why you think it happened and how you checked that, and the specific thing that you do differently now.

To build your own, use four parts: expectation, result, cause, change. Spend most of the time on cause and change. The stronger answers separate a bad process from bad luck. Sometimes you made the right decision on what you knew. Say so, and say how you would know.

\textbf{Illustration, with your own facts in the brackets.} ``On [project], I expected [result with a number], because [reason]. It came out at [number]. When I looked into it, the cause was [cause], which I checked by [check]. Part of that was bad luck, because [reason], but the part I own is [your error]. So I now [specific change], which I also used on [later example].''

\iqwrong{Choosing a disguised success (``we only beat the target by 5\%''), or blaming the team, the data or the market. Any of them leaves nothing for you to have learned.}
\iqfollow{What did you tell people when you saw it was going badly?}{Have a date in mind: when you knew, and when you said so. A long gap between those two is what they look for.}
\end{solution}

\begin{solution}{iq:iv2:behavioural-and-fit:16}
Choose a real mistake with a real cost, that you found or admitted yourself, and focus on what happened after you found it: whom you told, how quickly, what you did to fix it and the check that you added so that it cannot recur.

To build your own, pick something that is sized for your level and that you can tell in about two minutes. Avoid anything that hurts a former colleague, or breaches a confidence.

\textbf{Illustration, with your own facts in the brackets.} ``In [setting], I [mistake], and the cost was [number or effect]. I found it when [how]. Within [time], I told [person], and we [fix]. I now [check or habit], so the same thing would be caught by [how].''

\iqwrong{A mistake that is too small to count (``I once missed a comma''), or a mistake with no action afterwards. Worse is a long gap between finding it and telling someone. The error is forgivable, but the delay often is not.}
\iqfollow{What would have happened if you had said nothing?}{Answer truthfully, even if the answer is ``perhaps no one would ever have known''. Then you can explain why you told them anyway.}
\iqwhy{Desks can survive errors that are reported at once. What hurts firms is someone hiding a small loss until it is large. This is the answer an interviewer will remember.}
\end{solution}

\begin{solution}{iq:iv2:behavioural-and-fit:17}
Fix the facts first, then tell the person who decided, then correct the record, and then fix the cause. Do this today. You must tell them even if nobody ever would have noticed.

In order: first, check quickly that it really is wrong and how large it is. Keep that to minutes or an hour. Second, tell your manager and the person who used the number, directly, in the same day: what was wrong, the correct number and what it changes in their decision. Third, reissue the report with the correction marked. Fourth, find out why it happened, and add a test or a check. Fifth, check whether the same error is in anything else you sent.

\textbf{Illustration of the words.} ``The figure of [wrong number] on page [page] of last week's report is wrong. The right figure is [right number]. It came from my error in [step]. For the decision on [decision], this means [effect]. I am correcting the report now, and I am checking [what else].''

\iqwrong{Waiting until you have a complete analysis, or a good moment, or quietly correcting it next time. That turns a small error into a question about you.}
\iqfollow{And if your manager says to leave it, since it does not change much?}{Say why you would rather correct it, and put it in writing to them. If it affects a client or a regulatory figure, then escalate it through the proper route.}
\end{solution}

\begin{solution}{iq:iv2:behavioural-and-fit:18}
Tell a story in which you raised the issue with evidence, in private, and listened, and in which the end result was settled by facts or by a clear decision. Say plainly who turned out to be right, because you were asked. How you behaved matters more than the score, and so does going along with the decision.

To build your own, use five parts: what you each thought, what you brought, how it was settled, who turned out to be right, and what you learned about working with that person. A story in which you were wrong and said so can be stronger than one in which you won.

\textbf{Illustration, with your own facts in the brackets.} ``[Senior person] wanted [approach]. I thought [other approach] because [evidence]. I asked for [time], and I showed them [test or data]. They [outcome]. Afterwards [what happened, and who was right]. I learned [true lesson].''

\iqwrong{A story that ends with ``and I was proved right,'' with a fool on the other side. The other weak answer is ``I have never disagreed with anyone,'' which tells them that you won't speak up when a position is wrong.}
\end{solution}

\begin{solution}{iq:iv2:behavioural-and-fit:19}
Pick a case in which something looked too good or simply odd, and go through what made you doubt it, how you tested it and what the answer was. Finding a bug in your own good result is the strongest version.

To build your own, use: the signal that something was off, your list of possible causes, the cheapest test for each, and the conclusion. Good checks include a known answer on small data, a change in the time period, removing the top contributors, shuffling the labels, and an independent reimplementation.

\textbf{Illustration, with your own facts in the brackets.} ``[Model or analysis] gave [result], which seemed [too good or strange] because [benchmark]. I suspected [cause one] or [cause two]. I tested by [check], which showed [finding]. The true figure was [number]. I told [person] and we [action].''

\iqwrong{Choosing a case where you doubted a result only because you disliked it. Doubt needs a reason you can state, and an equal readiness to doubt a result you like. The other error is just rerunning the model and accepting it.}
\end{solution}

\begin{solution}{iq:iv2:behavioural-and-fit:20}
Talk to them first, in private, soon, and with specifics. Offer help, and agree on what they will deliver and by when. If that fails, adjust the plan with the team. Tell someone above you as soon as the deadline is really at risk, without surprises.

So the order is: first, find out why. They may be stuck, ill, or unclear on what is expected. Second, make the work small and visible, with a date for each part. Third, if they miss that, tell them plainly that you need to raise it, and then raise it. State facts about the work, and not about the person. In the meantime, protect the deadline by doing the most important parts, and say what will not get done.

\textbf{Illustration of the words.} ``I've noticed [part] isn't done, and we need it by [date]. Is something in the way? Could you have [small part] by [date]? If that is hard, tell me now, and we can change the plan.''

\iqwrong{Silently doing their part and resenting it, or going straight to the supervisor. The first hides a problem from the person who needs to know. The second ends the working relationship before you have even asked.}
\end{solution}

\begin{solution}{iq:iv2:behavioural-and-fit:21}
Do not ship it untested on your own authority, and do not simply refuse. Find out what is at stake, offer the fastest safe path, and make the person with the authority decide if the risk is to be taken.

The steps: first, ask what the change is, what it would earn and what happens if it is not in today. Second, size the risk: what is the worst it can do, and what limits or switches will stop it? Third, offer options, such as a smaller change, running it in shadow mode, only on a small position or a single symbol, behind a flag, or with a size limit, and tests limited to what the hour allows. Fourth, if the firm's process has a way to approve an emergency change, use it, and get the approval in writing, from whoever owns the risk.

\textbf{Illustration of the words.} ``I can't put that in blind, because if it goes wrong the loss is [worst case]. I can have [smaller safe version] ready in [time], with a limit of [size]. If you want the full version today, then [named owner] needs to approve it, and I'll prepare it so that we can pull it out in [time].''

\iqwrong{``Sure,'' because the trader is senior and in a hurry. Equally wrong is a flat no, with no alternative. A rushed release can lose a great deal of money very quickly.}
\iqfollow{What if it does go wrong?}{You switch it off first, then tell everyone, and only then look for the cause. Prepare the way to switch it off before you release it.}
\iqwhy{An engineer here is judged both on helping the desk and on not taking the firm down. This is the moment when those two meet.}
\end{solution}

\begin{solution}{iq:iv2:behavioural-and-fit:22}
Choose a real failure you owned. Tell it in time order: how you noticed, what you did to stop the harm, whom you told, how you found the cause, and what you changed so that it would be caught earlier. Stop the harm first, and look for the cause second.

To build your own, write a short post-mortem\index{post-mortem} as if it were for your team: impact, timeline, cause, what went well, and the actions taken. Give a number for the impact and for the time to detect and to recover. Be exact about your own part.

\textbf{Illustration, with your own facts in the brackets.} ``On [date], [system] did [failure]. We found out through [alert or person] after [time]. In the first ten minutes, I [action to stop the harm] and told [person]. The cause turned out to be [cause]. The impact was [number]. Afterwards, we added [test, alert or limit], so it would now be caught in [time].''

\iqwrong{Telling a story about a heroic all-nighter, when the point is that the failure should never have needed one. Blaming someone else for it is worse. If you have never had a failure, you have probably not owned much.}
\end{solution}

\begin{solution}{iq:iv2:behavioural-and-fit:23}
Choose a real deadline where something had to be dropped. Show how you decided: what you ranked as essential, whom you told what would not be done, and how the result came out. The decision is the interesting part, and not the long hours.

To build your own, use: what the deadline was and why it could not move, the options, how you chose, whom you told, and the outcome. Include what you got wrong in your first estimate. Keep your hours out of it, because long hours are not evidence of anything.

\textbf{Illustration, with your own facts in the brackets.} ``We had [time] to deliver [thing] because [reason]. At [point], I saw that we couldn't do all of it. I ranked [parts] by [criterion], dropped [part], and told [person] that morning. We delivered [result], and we finished [part] by [date]. Next time, I would [what you would do differently].''

\iqwrong{``I simply worked through the night, and did everything.'' It avoids the question, and it suggests that you may not tell people when you are in trouble.}
\end{solution}

\begin{solution}{iq:iv2:behavioural-and-fit:24}
For most people, take the 1 million, even though the coin toss averages 1.5 million. With logarithmic utility, you are indifferent at a starting wealth of exactly 1 million. Below that you take the certain amount, and above it you take the toss.

With wealth $W$ in millions, indifference means $\ln(W + 1) = \tfrac12\ln(W + 3) + \tfrac12\ln W$, so $(W + 1)^2 = W(W + 3)$, which gives $W = 1$. So someone with 100\,000 dollars should take the certain amount, and someone with 10 million should take the toss. A person's answer should depend on their wealth, and you should say so.

So the premium you give up, 0.5 million in expected value, pays for not risking a large part of your wealth. At $W = 10$, the certainty equivalent\index{certainty equivalent} of the toss is about 1.40 million, closer to its average. At $W = 0.1$, it is about 0.46 million.

\iqwrong{``I'd take the toss, because the expected value is higher,'' with no thought about wealth. An interviewer may follow this by asking you to bet your own house. The opposite error is treating a small bet the same way, for instance turning down a small, favourable bet because it is a bet.}
\iqfollow{What should a firm with 1 billion do?}{Take the toss, and every similar bet, because it is small relative to its capital. That is why a firm can take risk that a person could not.}
\iqwhy{It shows whether you separate expected value from what you can afford to lose, the idea behind all sizing.}
\end{solution}

\begin{solution}{iq:iv2:behavioural-and-fit:25}
Pick a real risk, not a financial boast. Say what you thought the odds and the costs were at the time, what you did to limit the downside, and judge the decision apart from the result. Say whether you would do it again on what you knew then.

To build your own, a move, a change of field, leaving a safe job, starting something, or betting on a skill all count. It should have cost you something real if it had failed. State your estimate then, even if it was rough. Then give the most interesting part: what you learned about how good that estimate was.

\textbf{Illustration, with your own facts in the brackets.} ``I [decision] in [year]. I thought I had about [chance] of [good result], and the downside was [cost], which I could live with because [reason]. It turned out [result]. I would [do or not do] it again on what I knew, because [reason]. The thing I misjudged was [what].''

\iqwrong{Saying that a decision was good because it worked, or describing something reckless as bravery. Someone who takes a large risk with no estimate of the odds is exactly what a firm does not want.}
\end{solution}

\begin{solution}{iq:iv2:behavioural-and-fit:26}
Both average 0.1 a day, and $B$ is the better strategy: its standard deviation is about half of $A$'s, so its ratio of average to standard deviation is about twice as good. Which is harder to live with depends on you. $A$ has rare large losses, and $B$ has long dull losing runs.

\begin{center}\small
\begin{tabular}{@{}lrr@{}}
 & $A$ & $B$ \\ \hline
Average per day & 0.10 & 0.10 \\
Standard deviation per day & 2.70 & 1.35 \\
Average over 250 days & 25 & 25 \\
Standard deviation over 250 days & 42.7 & 21.3 \\
Average over standard deviation & 0.037 & 0.074 \\
Annual average over standard deviation, 250 days & 0.59 & 1.17 \\
Typical longest run of losing days in 250 & 2 & 9 \\
Worst loss on one day & 8 & 1
\end{tabular}
\end{center}
With $A$, you have a good day 90\% of the time, and then a loss of 8 removes eight good days. With $B$, you lose on most days, and a typical year holds a run of about 9 losses in a row, which is why many people drop it at the worst moment. The run lengths come from an exact count of the median longest run over 250 days.

A sound answer says that $B$ is better on the numbers, that $A$ looks better until you see the tail, and how you know which you could keep running. Before answering, ask whether the tail of $A$ is really only 8.

\iqwrong{``$A$, because it wins 90\% of the time.'' How often you win says nothing by itself, and a high hit rate is often how a seller of tail risk looks.}
\iqfollow{What would you check before trusting $A$'s numbers?}{How long the record is. A 10\% loss rate is seen quickly, but a much rarer and larger loss may not have shown up yet.}
\iqwhy{Your temperament has to fit the strategy's pattern of results. People abandon good strategies during the normal bad runs.}
\end{solution}

\begin{solution}{iq:iv2:behavioural-and-fit:27}
Stop adding risk, reduce the position as the limit requires, tell your manager at once with the facts, and then spend the day finding out whether the loss came from your decisions or from something that broke. Do not try to make it back.

The order: first, follow the rule. A loss limit exists for this moment, so flatten or reduce as it says. Second, check that the loss is real: are the marks, positions and data right? Third, tell your manager in numbers: how much, from what, what you hold now and what you plan. Fourth, separate the loss into parts, for example the market's move, your position size and your execution. Fifth, spend what remains of the day on work you can do well, such as reviewing the trade, or watching without trading.

\textbf{Illustration of the words.} ``I'm down [amount], which is my limit. It came mainly from [position], because [cause]. I have [reduced to size]. The marks check out. I think [judgement]. I will not add today unless you want me to.''

For the next day, return at a reduced size, and raise it only when your process is working again. Ask whether the original idea still stands, as a new trade.

\iqwrong{Trading harder to recover the loss. It is the classic way to turn a bad day into a far worse one. Hiding the size from your manager is worse.}
\iqfollow{What if you are sure that the position will come back?}{Everyone always is. If you were right about the position, the desk can put it on again with a clear head. A limit that you override when you are confident is no limit.}
\iqwhy{Bad days happen to every trader. What the firm cannot accept is a person whose behaviour gets worse after a loss.}
\end{solution}

\begin{solution}{iq:iv2:behavioural-and-fit:28}
Treat it as a breach, not a triumph. Tell your manager yourself, work out what the rules would have made and what the extra risk was, and decide openly whether the rule or your behaviour should change. A broken rule that paid off is more dangerous than one that lost.

The profit does not show that the decision was good. Separate the pieces: what would the rule-size trade have made, and what did the extra size add? What would the extra size have lost, if the event had gone the other way? If that loss would have been more than you could accept, then you took a risk that you had decided not to take.

Next, ask why. Either the rule is wrong, in which case you change it deliberately, in writing, with evidence, when you are calm. Or the rule is right, and you broke it because you felt sure, which is the feeling to watch for. Next time, write the reason for going over size before you trade, and ask someone else to look at it.

\textbf{Illustration of the words.} ``My best trade this year broke my sizing rule. Here is what the rule would have made, and here is what I would have lost if it had gone the other way. I want to talk about whether the rule should change. It should not change just because this one worked.''

\iqwrong{Quietly keeping the profit, or congratulating yourself on your instinct. A manager who finds out later will wonder what else went unreported.}
\iqfollow{What if your manager is delighted, and tells you to do more of it?}{Show them the same numbers, and ask for an explicit change to your limits if that is what they want. You should not leave it as a feeling.}
\end{solution}

\begin{solution}{iq:iv2:behavioural-and-fit:29}
Check your facts quietly and quickly, then raise it through the proper route: first your manager, unless your manager is involved, and then risk, product control or compliance. Do not ignore it. Do not accuse anyone before you have facts, and do not confront them alone with an accusation. If nothing happens after that, take independent legal advice on going to the regulator.

First, what do you actually know? Compare the marks with independent sources, broker quotes or the firm's independent price checks, and look at how the difference has changed over time. An illiquid position may have an honest reason for the difference, so write down what would explain it.

Next, ask a neutral question if that is natural: ``I can't get to this mark. What am I missing?'' It often settles it. If it does not, take your written facts to your manager. If your manager is involved, or nothing changes, use the firm's own channel: risk, product control, compliance or a whistleblowing line. Keep a dated record of what you did. You need not be certain in order to report it. Establishing the truth is their job.

If you have used the firm's channels and still nothing happens, take independent legal advice on reporting to the regulator. Legal protection for whistleblowers varies from country to country. Do not take documents out of the firm, or go to anyone else outside it, before you have that advice.

\iqwrong{``It's not my position, so it's not my business,'' or fearing for your bonus. The first can make you part of it if it is true. Mismarking is how a hidden loss can grow for years. The opposite error is an open accusation on a desk, on the basis of one price.}
\iqfollow{And if reporting it would hurt your career?}{It may, and you should say that you know that. Say also that staying silent, on something that might be serious, is the larger risk to you and to the firm.}
\end{solution}

\begin{solution}{iq:iv2:behavioural-and-fit:30}
Stay calm, check once out loud, and if you still believe it, say so politely, and show your reasoning so they can point to the step they disagree with. If they are right, say so at once, name your mistake and carry on.

Often the interviewer is testing you. Either they want to see if you hold a correct answer under pressure, or they have misheard, or they are thinking of a different version of the problem. So the useful first move is to ask: ``Which step do you disagree with?'' Then check a simple case, or try another method. Say what you find.

\textbf{Illustration of the words.} ``Let me check that. I assumed [assumption], which gives [step]. A small case also gives [result]. So I still get [answer]. Have I misread the question, or do you see an error in a step?'' And if you find that they are right: ``You're right. My error was [step]. That changes the answer to [answer].''

Both of these are good outcomes, and the second is not a failure. Abandoning at once a correct answer you can defend is a poor result. So is arguing for one you have just found to be wrong.

\iqwrong{Folding at the first frown, which suggests your view will not survive a trader's pressure. Getting defensive or arguing without showing any working is just as bad.}
\iqfollow{What if they insist, and you are sure?}{Say that you see it differently, and that you are happy to send a short proof afterwards. Move on, and do not let it spoil the rest of the interview.}
\end{solution}
